% Gliederung Stand 30.10.2026. Stand der Technik, ~8 Seiten.
\chapter{Fähigkeiten humanoider Roboter und verwandte Arbeiten}
\label{chap:faehigkeiten}
% Kurze Kapiteleinleitung vor 4.1.

\section{Fortbewegung}
\label{sec:fortbewegung_sdt}

\subsection{Modellbasierte Gleichgewichtsregelung}
\label{sec:gleichgewicht}
% Schwerpunkt, Druckpunkt, PD-Regelung. Kandidaten: Saeedvand2019 (COM/COP, ZMP), Tong2024 (Grenzen ZMP), Yu2026 (PD-Schwerpunktausgleich).

\subsection{Lernbasierte Lauf-Regler}
\label{sec:rl_laufen}
% Reinforcement Learning. Kandidaten: Radosavovic2024, Tong2024 (RL und Sim-Real-Unterschied).

\section{Manipulation}
\label{sec:manipulation_sdt}

\subsection{Inverse Kinematik redundanter Arme}
\label{sec:ik}
% Gedämpfte Pseudoinverse, Nullraum. Kandidaten: Kucuk2006 (Grundbegriffe FK/IK, DH, Singularität), Aristidou2009 (Pseudoinverse, DLS, Nullraum), Aristidou2018 (Redundanz, Lösungsfamilien, DLS, Nullraum).

\subsection{Bahnplanung}
\label{sec:bahnplanung_sdt}
% Sampling-basiert (RRT-Connect), optimierungsbasiert (CHOMP), Frameworks (MoveIt, OMPL). Kandidat: Yu2026 (Polynom 5. Grades).

\subsection{Kollisionsmodelle}
\label{sec:kollision_sdt}
% Kugeln, konvexe Hüllen.

\section{Bedienung und Teleoperation humanoider Roboter}
\label{sec:teleoperation_sdt}
% Kandidaten: Saeedvand2019 (autonom/teilautonom/ferngesteuert, Fernsteuerung nötig), Tong2024 (Trends Teleoperation).

\section{Verwandte Arbeiten}
\label{sec:verwandte_arbeiten}
% Bestehende Umgebungen und Software-Stacks für den G1 und vergleichbare Humanoide
% (u. a. HuCeBot g1pilot, Unitree-Beispiele, unitree_rl_mjlab, WBC-AGILE). Kandidat: Yu2026 (G1 nur in MATLAB/Simulink).
